% GENERATED FILE. Edit canonical lessons, then run npm run generate:books.
% canonical-source-hash: fnv1a64:8ee5d84bd43a688e
% canonical-lessons: GE-C486-muedigkeit, GE-C486-erschoepft, GE-C486-gaehnen, GE-C486-schnarchen, GE-C486-albtraum

\chapter{Müde und erschöpft}
\label{ch:ge-mude-und-erschopft}
\hlchaptermodality{486}

\begin{chapteropening}
\textbf{By the end of this chapter:} \emph{I can talk about tiredness and sleep problems.}

Complete the last lesson of chapter 486: I can talk about tiredness and sleep problems.
\end{chapteropening}

\section[die Müdigkeit]{die Müdigkeit — tiredness}

\label{lesson:GE-C486-muedigkeit}



\begin{quote}
Before the new one: say the German for ankle, then the German for hip.
\end{quote}

\subsection*{You'll want to know: die Müdigkeit}
\textbf{die Müdigkeit} — \textquotedblleft{}tiredness\textquotedblright{}.

\begin{quote}
Nach dem Mittagessen kommt bei mir oft die Müdigkeit.
\end{quote}

\begin{quote}
\emph{After lunch I often get tired.}
\end{quote}

No plural. Built on \emph{müde} — tired.

\begin{cousinweb}
\emph{Müdigkeit} is built on \emph{müde} \textquotedblleft{}tired\textquotedblright{} with \emph{-keit}; \emph{müde} is inherited from Old High German \emph{muodi}, from Proto-Germanic \emph{mōþijaz} (EtymDB).
\end{cousinweb}

\begin{grammarlens}[title={What you've built}]
One more word for telling, explaining and describing at B1.
\end{grammarlens}

\subsection*{Your turn}
\begin{itemize}
\raggedright
  \item \emph{Say it:} \emph{die Müdigkeit}
  \item \emph{Say it:} \emph{die Müdigkeit}, once more
  \item \emph{Say it:} a sentence of your own with \emph{die Müdigkeit}
  \item \emph{Recall:} say the German for ankle, then the German for hip, then say \emph{die Müdigkeit} again
\end{itemize}

\begin{tcolorbox}[breakable,colback=teal!4,colframe=teal!35!black,title={Before you move on}]
What does die Müdigkeit mean? (\textquotedblleft{}Tiredness\textquotedblright{}.) Say it once more.
\end{tcolorbox}
\section[erschöpft]{erschöpft — exhausted}

\label{lesson:GE-C486-erschoepft}



\begin{quote}
Before the new one: say the German for hip, then the German for tiredness.
\end{quote}

\subsection*{You'll want to know: erschöpft}
\textbf{erschöpft} — \textquotedblleft{}exhausted\textquotedblright{}.

\begin{quote}
Nach dem Umzug waren wir völlig erschöpft.
\end{quote}

\begin{quote}
\emph{After the move we were completely exhausted.}
\end{quote}

The participle of \emph{erschöpfen} — to exhaust. The noun: \emph{die Erschöpfung}.

\begin{cousinweb}
\emph{erschöpft} is built on \emph{er-} and \emph{schöpfen} \textquotedblleft{}to scoop, to draw water\textquotedblright{}, from Middle High German \emph{schepfen} and Old High German \emph{scaphan} (EtymDB): scooped empty.
\end{cousinweb}

\begin{grammarlens}[title={What you've built}]
One more word for telling, explaining and describing at B1.
\end{grammarlens}

\subsection*{Your turn}
\begin{itemize}
\raggedright
  \item \emph{Say it:} \emph{erschöpft}
  \item \emph{Say it:} \emph{erschöpft}, once more
  \item \emph{Say it:} a sentence of your own with \emph{erschöpft}
  \item \emph{Recall:} say the German for hip, then the German for tiredness, then say \emph{erschöpft} again
\end{itemize}

\begin{tcolorbox}[breakable,colback=teal!4,colframe=teal!35!black,title={Before you move on}]
What does erschöpft mean? (\textquotedblleft{}Exhausted\textquotedblright{}.) Say it once more.
\end{tcolorbox}
\section[gähnen]{gähnen — to yawn}

\label{lesson:GE-C486-gaehnen}



\begin{quote}
Before the new one: say the German for tiredness, then the German for exhausted.
\end{quote}

\subsection*{You'll want to know: gähnen}
\textbf{gähnen} — \textquotedblleft{}to yawn\textquotedblright{}.

\begin{quote}
In der langen Besprechung gähnten alle.
\end{quote}

\begin{quote}
\emph{Everyone was yawning in the long meeting.}
\end{quote}

Regular: \emph{hat gegähnt}. \emph{Gähnende Leere} — total emptiness.

\begin{cousinweb}
\emph{gähnen} is inherited from Old High German \emph{ginēn} \textquotedblleft{}to be wide open, to gape\textquotedblright{} (EtymDB).
\end{cousinweb}

\begin{grammarlens}[title={What you've built}]
One more word for telling, explaining and describing at B1.
\end{grammarlens}

\subsection*{Your turn}
\begin{itemize}
\raggedright
  \item \emph{Say it:} \emph{gähnen}
  \item \emph{Say it:} \emph{gähnen}, once more
  \item \emph{Say it:} a sentence of your own with \emph{gähnen}
  \item \emph{Recall:} say the German for tiredness, then the German for exhausted, then say \emph{gähnen} again
\end{itemize}

\begin{tcolorbox}[breakable,colback=teal!4,colframe=teal!35!black,title={Before you move on}]
What does gähnen mean? (\textquotedblleft{}To yawn\textquotedblright{}.) Say it once more.
\end{tcolorbox}
\section[schnarchen]{schnarchen — to snore}

\label{lesson:GE-C486-schnarchen}



\begin{quote}
Before the new one: say the German for exhausted, then the German for to yawn.
\end{quote}

\subsection*{You'll want to know: schnarchen}
\textbf{schnarchen} — \textquotedblleft{}to snore\textquotedblright{}.

\begin{quote}
Mein Mann schnarcht so laut, dass ich nicht schlafen kann.
\end{quote}

\begin{quote}
\emph{My husband snores so loudly that I can't sleep.}
\end{quote}

Regular: \emph{hat geschnarcht}.

\begin{cousinweb}
\emph{schnarchen} is inherited from Middle High German \emph{snarchen} (EtymDB), probably an imitative word; English \emph{snore} and \emph{snort} are of the same kind.
\end{cousinweb}

\begin{grammarlens}[title={What you've built}]
One more word for telling, explaining and describing at B1.
\end{grammarlens}

\subsection*{Your turn}
\begin{itemize}
\raggedright
  \item \emph{Say it:} \emph{schnarchen}
  \item \emph{Say it:} \emph{schnarchen}, once more
  \item \emph{Say it:} a sentence of your own with \emph{schnarchen}
  \item \emph{Recall:} say the German for exhausted, then the German for to yawn, then say \emph{schnarchen} again
\end{itemize}

\begin{tcolorbox}[breakable,colback=teal!4,colframe=teal!35!black,title={Before you move on}]
What does schnarchen mean? (\textquotedblleft{}To snore\textquotedblright{}.) Say it once more.
\end{tcolorbox}
\section[der Albtraum]{der Albtraum — nightmare}

\label{lesson:GE-C486-albtraum}



\begin{quote}
Before the new one: say the German for to yawn, then the German for to snore.
\end{quote}

\subsection*{You'll want to know: der Albtraum}
\textbf{der Albtraum} — \textquotedblleft{}nightmare\textquotedblright{}.

\begin{quote}
Nach dem Horrorfilm hatte ich einen Albtraum.
\end{quote}

\begin{quote}
\emph{After the horror film I had a nightmare.}
\end{quote}

Plural: \emph{die Albträume}. Also spelt \emph{der Alptraum}.

\begin{cousinweb}
\emph{Albtraum} is built on \emph{Alb} \textquotedblleft{}elf, night spirit\textquotedblright{}, inherited from Middle High German \emph{alb} and Old High German \emph{alp}, from Proto-Germanic \emph{albiz} (EtymDB), and \emph{Traum} \textquotedblleft{}dream\textquotedblright{}.
\end{cousinweb}

\begin{grammarlens}[title={What you've built}]
One more word for telling, explaining and describing at B1.
\end{grammarlens}

\subsection*{Your turn}
\begin{itemize}
\raggedright
  \item \emph{Say it:} \emph{der Albtraum}
  \item \emph{Say it:} \emph{der Albtraum}, once more
  \item \emph{Say it:} a sentence of your own with \emph{der Albtraum}
  \item \emph{Recall:} say the German for to yawn, then the German for to snore, then say \emph{der Albtraum} again
\end{itemize}

\begin{tcolorbox}[breakable,colback=teal!4,colframe=teal!35!black,title={Before you move on}]
What does der Albtraum mean? (\textquotedblleft{}Nightmare\textquotedblright{}.) Say it once more.
\end{tcolorbox}
